\chapter{الاحتمالات على الفضاءات القابلة للعد}\label{ch:b2:proba}

تطور الفصول الثلاثة الأخيرة نظرية الاحتمالات الحديثة: القياسات
الاحتمالية على الفضاءات العينية \omterm{def:b2:structures:countable}{القابلة للعد}، والمتغيرات العشوائية
المتقطعة، والدوال المولّدة. فتكتسب النظرية المنتهية في مجلد الثانوية
بنيتها التحتية الكاملة: إذ تحل الجمعية $\sigma$ محل الجمعية المنتهية،
وتكون آلة العائلات \omterm{def:b2:series:summable}{القابلة للجمع} في \cref{ch:b2:series} هي بالضبط ما يجعل الفضاءات
العينية اللانهائية قابلة للعمل. والنتيجتان المحوريتان هنا هما اتصال
الاحتمال على المتتاليات الرتيبة من \omterm{def:b2:proba:space}{الأحداث} ومبرهنة بوريل--كانتيلي.

\section{الفضاءات الاحتمالية}

\begin{definition}[الفضاء الاحتمالي القابل للعد]\label{def:b2:proba:space}
لتكن $\Omega$ مجموعة غير خالية منتهية أو \omterm{def:b2:structures:countable}{قابلة للعد} (وهي \emph{الفضاء
العيني}\index{فضاء عيني}). \emph{القياس الاحتمالي}\index{قياس احتمالي}
على $\Omega$ هو تطبيق $\P$ من مجموعة $\mathcal{P}(\Omega)$ جميع أجزاء $\Omega$
(\emph{الأحداث}\index{حدث}) نحو $[0, 1]$ يحقق:
\begin{enumerate}
\item $\P(\Omega) = 1$؛
\item (الجمعية $\sigma$) من أجل كل متتالية $(A_n)_{n\in\N}$ من الأحداث المنفصلة
مثنى مثنى،
\[
\P\Bigl(\,\bigcup_{n \in \N} A_n\Bigr)
= \sum_{n=0}^{\infty} \P(A_n) .
\]
\end{enumerate}
والزوج $(\Omega, \P)$ هو \emph{فضاء احتمالي}\index{فضاء احتمالي} (قابل للعد).
\end{definition}

\begin{remark}
على مجموعة $\Omega$ \omterm{def:b2:structures:countable}{قابلة للعد} يمكن أن نأخذ جميع الأجزاء أحداثا؛ أما
على الفضاءات غير \omterm{def:b2:structures:countable}{القابلة للعد} (كما تقتضيه النماذج المتصلة في السنة
الثالثة) فلم يعد ذلك ممكنا، ونقصر $\P$ على مجموعة مناسبة من
\omterm{def:b2:proba:space}{الأحداث}، هي \emph{جبر $\sigma$}. وتصمد جميع صيغ هذا الفصل أمام ذلك
التعميم حرفيا.
\end{remark}

\begin{proposition}[القواعد الأولية]\label{prop:b2:proba:rules}
من أجل أحداث $A, B$ \omterm{def:b2:proba:space}{وقياس احتمالي} $\P$: $\P(\emptyset) = 0$؛ والتطبيق $\P$
جمعي منتهيا؛ و$\P(A^c) = 1 -
\P(A)$؛ وإذا كان $A \subseteq B$ فإن $\P(A) \leq \P(B)$؛ و
\[
\P(A \cup B) = \P(A) + \P(B) - \P(A \cap B) .
\]
\end{proposition}

\begin{proof}
تطبيق الجمعية $\sigma$ على $A_0 = \Omega$، $A_n = \emptyset$ ($n \geq 1$) يعطي $1 = 1 + \sum_{n\geq1}\P(\emptyset)$، ومنه
$\P(\emptyset) = 0$؛ ثم إن حشو اتحاد منفصل منته بمجموعات خالية يعطي الجمعية
المنتهية. أما الباقي فينتج كما في الحالة المنتهية (مجلد الثانوية):
$1 = \P(A) + \P(A^c)$ انطلاقا من $\Omega = A \sqcup A^c$؛ و$\P(B) = \P(A) + \P(B \setminus A)
\geq \P(A)$ عندما $A \subseteq B$؛ وبالتفكيك إلى ثلاث
قطع منفصلة،
\begin{align*}
\P(A \cup B) &= \P(A \setminus B) + \P(B \setminus A) +
\P(A \cap B)\\
&= \bigl(\P(A) - \P(A\cap B)\bigr) + \bigl(\P(B) - \P(A\cap
B)\bigr) + \P(A \cap B),
\end{align*}
وهو الاحتواء والاستبعاد؛ والصيغة العامة ذات $n$ مجموعة هي \cref{exo:b2:proba:4}.
\end{proof}

\begin{proposition}[التوزيعات على فضاء قابل للعد]\label{prop:b2:proba:distribution}
إعطاء \omterm{def:b2:proba:space}{قياس احتمالي} على مجموعة $\Omega =
\{\omega_0, \omega_1, \dots\}$ \omterm{def:b2:structures:countable}{قابلة للعد} يعادل بالضبط إعطاء
أوزان $p_i = \P(\{\omega_i\}) \geq 0$ تحقق $\sum_i p_i = 1$؛ وعندئذ يكون، من أجل كل $A \subseteq \Omega$،
\[
\P(A) = \sum_{\omega \in A} \P(\{\omega\}) ,
\]
وهو مجموع جزئي (متقارب مطلقا) للعائلة $(p_i)$.
\end{proposition}

\begin{proof}
بمعطى $\P$، تشكل المجموعات الأحادية $\{\omega\}$، $\omega \in A$، تغطية منفصلة
\omterm{def:b2:structures:countable}{قابلة للعد} للمجموعة $A$، ومنه تفرض الجمعية $\sigma$
\[
\P(A) = \sum_{\omega\in A}\P(\{\omega\}),
\]
وهو مجموع جزئي غير مشروط \omterm{def:b2:series:summable}{للعائلة القابلة للجمع} الموجبة $(p_i)$ ---
فإعادة الترتيب غير ضارة بالضبط لأن الحدود موجبة (\cref{ch:b2:series})؛ وبوجه خاص
$\sum_ip_i = \P(\Omega) = 1$.
وبالعكس، بمعطى أوزان موجبة مجموعها الكلي $1$، نعرّف $\P(A) = \sum_{\omega \in A}p_\omega$: \omterm{def:b2:series:summable}{فالعائلة
قابلة للجمع}، والجمعية $\sigma$ هي بالضبط مبرهنة الجمع بالرزم في \cref{ch:b2:series}
مطبقة على تجزئة $\bigcup A_n$ إلى المجموعات $A_n$.
\end{proof}

\begin{example}[النموذج الهندسي: انتظار أول صورة]\label{ex:b2:proba:geometric}
نرمي قطعة نقود احتمال ظهور الصورة فيها $p \in \intoo{0}{1}$ مرارا وتكرارا، ونجعل
$\Omega = \N^* \cup \{\infty\}$ يسجل رتبة أول صورة. والأوزان الطبيعية هي
\[
\P(\{k\}) = (1 - p)^{k-1}p
\quad (k \in \N^*),
\qquad
\P(\{\infty\}) = 0 ,
\]
وهي \omterm{def:b2:proba:space}{قياس احتمالي} لأن $\sum_{k\geq1}(1-p)^{k-1}p =
\frac{p}{1 - (1-p)} = 1$: فباحتمال $1$ تنتهي اللعبة --- لكن
يجب أن يظل \omterm{def:b2:proba:space}{الفضاء العيني} حاويا لإمكان ألا تنتهي. والجمعية \omterm{def:b2:structures:countable}{القابلة
للعد} هي ما يسمح لنا بأن نؤكد $\P(\text{تنتهي اللعبة}) = \sum_k \P(\{k\})$.
\end{example}

\begin{theorem}[الاتصال الرتيب]\label{thm:b2:proba:continuity}
لتكن $(A_n)$ متتالية من \omterm{def:b2:proba:space}{الأحداث}.
\begin{enumerate}
\item إذا كان $A_n \subseteq A_{n+1}$ من أجل كل $n$ (\emph{متزايدة})، فإن $\P\bigl(\bigcup_n A_n\bigr) = \lim_{n\to\infty} \P(A_n)$.
\item إذا كان $A_n \supseteq A_{n+1}$ من أجل كل $n$ (\emph{متناقصة})، فإن $\P\bigl(\bigcap_n A_n\bigr) = \lim_{n\to\infty} \P(A_n)$.
\end{enumerate}
\end{theorem}

\begin{proof}
\emph{1.} نفصل \omterm{def:b2:proba:space}{الأحداث}: نضع $B_0 = A_0$ و$B_n = A_n \setminus
A_{n-1}$. \omterm{def:b2:proba:space}{فالأحداث} $B_n$ منفصلة
مثنى مثنى مع $\bigcup_{k \leq n}
B_k = A_n$ و$\bigcup_n B_n = \bigcup_n A_n$. وبالجمعية $\sigma$ والجمعية المنتهية،
\[
\P\Bigl(\bigcup_n A_n\Bigr)
= \sum_{n=0}^\infty \P(B_n)
= \lim_{N\to\infty}\sum_{n=0}^N \P(B_n)
= \lim_{N\to\infty}\P(A_N) .
\]
\emph{2.} نمر إلى المتممات: فالمتتالية $(A_n^c)$ متزايدة واتحادها $\bigl(\bigcap A_n\bigr)^c$،
ونطبق البند 1: $1 - \P(\bigcap A_n) = \lim (1 - \P(A_n))$.
\end{proof}

\begin{corollary}[الجمعية التحتية القابلة للعد]\label{cor:b2:proba:subadd}
من أجل أي متتالية من \omterm{def:b2:proba:space}{الأحداث}، $\P\bigl(\bigcup_n A_n\bigr) \leq
\sum_{n=0}^\infty \P(A_n)$.
\end{corollary}

\begin{proof}
تنتج الجمعية التحتية المنتهية $\P(A_0 \cup \dots \cup A_N) \leq
\sum_0^N \P(A_n)$ من الاحتواء والاستبعاد بالتراجع
(أو من الجمعية على \omterm{def:b2:proba:space}{الأحداث} المفصولة $B_n \subseteq A_n$). ولنجعل $N \to \infty$: يتقارب
الطرف الأيسر نحو $\P(\bigcup_n A_n)$ بالاتصال الرتيب مطبقا على المتتالية
المتزايدة $C_N = A_0 \cup \dots \cup A_N$.
\end{proof}

\begin{example}[حصر الاتحاد: فج لكنه لا يُهزم]\label{ex:b2:proba:unionbound}
الجمعية التحتية مع عدد منته من \omterm{def:b2:proba:space}{الأحداث} --- وهي \emph{حصر الاتحاد}
--- تبادل الدقة بالشمولية. ففي مسألة أعياد الميلاد مع $23$ شخصا،
يعطي حصر احتمال التصادم بالمجموع على الأزواج
\[
\P(\text{تصادم}) \leq \binom{23}2\cdot\frac1{365}
= \frac{253}{365} \approx 0.693 ,
\]
مقابل القيمة الحقيقية $0.507$: أي بفارق واسع، لأن التصادمات تتراكب.
ومع ذلك فإن الحصر لا يحتاج إلى \emph{أي} استقلال، ولا إلى قانون
مشترك، ولا إلى شيء سوى احتمالات الأزواج --- ولهذا يكون حصر الاتحاد،
في مسألة نهاية الأسبوع وفي كل \cref{ch:b2:randomvar}، أول أداة تُسلّ: فإن صادف أن
كان صغيرا، حُسم الأمر دون أي نمذجة إضافية.
\end{example}

\begin{example}[ستة تأتي في النهاية]\label{ex:b2:proba:sixeventually}
نرمي نردا متزنا إلى الأبد ونجعل $B_n = {}$``ظهور ستة واحدة على الأقل في
أول $n$ رمية''، وهي متتالية متزايدة من \omterm{def:b2:proba:space}{الأحداث} مع $\P(B_n) = 1 - (5/6)^n$. ويعطي
الاتصال الرتيب
\[
\P(\text{تظهر ستة في النهاية})
= \P\Bigl(\bigcup_nB_n\Bigr)
= \lim_n\bigl(1 - (5/6)^n\bigr) = 1 .
\]
والمهم ليس النهاية (البديهية) بل الخطوة المنطقية: فقولنا ``في
النهاية'' \omterm{def:b2:proba:space}{حدث} يتعلق \emph{بعدد لانهائي} من الرميات، خارج متناول
الجمعية المنتهية، والاتصال الرتيب --- أي الجمعية $\sigma$ --- هو
بالضبط البديهية التي تسند إليه احتمالا. وكل نص شبه أكيد في بقية هذا
الكتاب يمر عبر هذا الباب الضيق نفسه.
\end{example}

\section{الشرط والاستقلال}

\begin{definition}[الاحتمال الشرطي]\label{def:b2:proba:conditional}
من أجل حدثين $A, B$ مع $\P(B) > 0$، يكون \emph{الاحتمال الشرطي}\index{احتمال شرطي}
\omterm{def:b2:proba:space}{للحدث} $A$ علما \omterm{def:b2:proba:space}{بالحدث} $B$ هو
\[
\P(A \mid B) = \frac{\P(A \cap B)}{\P(B)} .
\]
والتطبيق $A \mapsto \P(A \mid B)$ هو نفسه \omterm{def:b2:proba:space}{قياس احتمالي} على $\Omega$.
\end{definition}

\begin{remark}
كون $A \mapsto \pcond BA$ قياسا احتماليا من جديد يستحق وقفة: فالمقدار $\pcond B\Omega = 1$
والجمعية $\sigma$ يمران عبر القسمة لأن التقاطع مع $B$ يحترم
الاتحادات المنفصلة. والنتيجة العملية: يمكن تطبيق كل متطابقة من هذا
الفصل --- الاحتواء والاستبعاد، والاتصال الرتيب، وبوريل--كانتيلي ---
\emph{بعد} الشرط، دون أي براهين جديدة. ويعمل أهل الاحتمالات باستمرار
``تحت شرط $\pcond B{\cdot}$'' لهذا السبب بالضبط.
\end{remark}

\begin{example}[الشرط قد يخلق انتظاما]\label{ex:b2:proba:sumseven}
نرمي نردين متزنين ونشترط أن يكون المجموع $7$: فمن أجل كل $k \in \intint16$،
\[
\pcond{\{S = 7\}}{X = k}
= \frac{\P(X = k,\ Y = 7 - k)}{\P(S = 7)}
= \frac{1/36}{6/36} = \frac16 :
\]
أي إن النرد الأول، علما بأن المجموع $7$، منتظم تماما --- فالمجموع
$7$ هو الوحيد المتوافق مع كل وجه، ومنه فإن الشرط يمحو كل معلومة
عن $X$. وأي مجموع آخر يشوّه القانون (فعلما بأن $S = 4$، يكون
النرد الأول منتظما على $\{1, 2, 3\}$ فقط). وحساب قانون شرطي يعني إعادة توحيد
الأوزان المشتركة على طول \omterm{def:b2:proba:space}{حدث} الشرط، لا أكثر.
\end{example}

\begin{example}[السحب الثاني بجودة الأول]\label{ex:b2:proba:seconddraw}
يحوي جرة $3$ كرة بيضاء و$2$ كرة سوداء؛ ونسحب كرتين دون
إرجاع. والكل متفق على أن $\P(W_1) = \frac35$؛ فما قيمة $\P(W_2)$؟ بالاحتمالات
الكلية على السحب الأول:
\[
\P(W_2) = \pcond{W_1}{W_2}\,\P(W_1) +
\pcond{B_1}{W_2}\,\P(B_1)
= \frac24\cdot\frac35 + \frac34\cdot\frac25
= \frac{12}{20} = \frac35 :
\]
أي $\P(W_1)$ بالضبط. ولم يكن أي حساب لازما: فبالتناظر، لكل كرة الاحتمال
نفسه في أن تكون الكرة الثانية المسحوبة، ومنه فإن للسحب الثاني ---
\emph{دون شرط} --- القانون نفسه للسحب الأول. فالشرط على نتيجة السحب
الأول يغير الحظوظ؛ أما عدم معرفتها فلا. وتعود حجة التبادلية هذه في
الفصل التالي من أجل المعاينة دون إرجاع، حيث تعطي المتوسط فوق الهندسي
$np$ دون أي متطابقات ثنائية الحد البتة.
\end{example}

\begin{theorem}[الاحتمالات المركبة، والاحتمالات الكلية، وبايز]\label{thm:b2:proba:bayes}
\leavevmode
\begin{enumerate}
\item (قاعدة السلسلة) إذا كان $\P(A_1 \cap \dots \cap A_{n-1}) > 0$، فإن
\[
\P(A_1 \cap \dots \cap A_n)
= \P(A_1)\,\P(A_2 \mid A_1)\cdots
\P(A_n \mid A_1 \cap \dots \cap A_{n-1}) .
\]
\item (الاحتمالات الكلية) إذا كانت $(B_i)_{i \in I}$ تجزئة منتهية أو \omterm{def:b2:structures:countable}{قابلة
للعد} للفضاء $\Omega$ مع $\P(B_i) > 0$، فإن من أجل كل \omterm{def:b2:proba:space}{حدث} $A$:
\[
\P(A) = \sum_{i \in I} \P(A \mid B_i)\,\P(B_i) .
\]
\item (بايز) تحت الفرضيات نفسها، وإذا كان علاوة على ذلك $\P(A) > 0$:
\[
\P(B_j \mid A)
= \frac{\P(A \mid B_j)\,\P(B_j)}
       {\sum_{i \in I} \P(A \mid B_i)\,\P(B_i)} .
\]
\end{enumerate}
\end{theorem}

\begin{proof}
\emph{1.} نكتب كل \omterm{def:b2:proba:conditional}{احتمال شرطي} على صورة خارج قسمة: فالطرف الأيمن هو
\[
\P(A_1)\cdot\frac{\P(A_1 \cap A_2)}{\P(A_1)}\cdot
\frac{\P(A_1 \cap A_2 \cap A_3)}{\P(A_1 \cap A_2)}\cdots
\frac{\P(A_1 \cap \dots \cap A_n)}{\P(A_1 \cap \dots \cap
A_{n-1})},
\]
وهو جداء تلسكوبي: فكل مقام يبسّط البسط السابق، ولا يبقى إلا $\P(A_1 \cap \dots \cap A_n)$.
وجميع المقامات $\geq \P(A_1 \cap \dots \cap A_{n-1}) >
0$ بالرتابة، ومنه لا ينعدم شيء. (وهذا بالضبط ما
تحرسه الفرضية: فالشرط على \omterm{def:b2:proba:space}{حدث} احتماله معدوم غير معرَّف.)
\emph{2.} المجموعات $A \cap B_i$ منفصلة مثنى مثنى واتحادها $A$؛ ونطبق
الجمعية ($\sigma$) وتعريف الشرط.
\emph{3.} طرفا $\P(B_j \mid A)\P(A) = \P(A \mid
B_j)\P(B_j)$ يساويان $\P(A \cap B_j)$؛ ونقسم على $\P(A)$ وننشر
$\P(A)$ بالاحتمالات الكلية.
\end{proof}

\begin{example}[تصادم أعياد الميلاد، بقاعدة السلسلة]\label{ex:b2:proba:birthday}
مع $n$ شخصا أعياد ميلادهم مستقلة ومنتظمة على $365$ يوما، نجعل
$D_n = {}$``جميع أعياد الميلاد $n$ مختلفة''. وبالشرط شخصا بعد شخص
(قاعدة السلسلة):
\[
\P(D_n) = \prod_{k=1}^{n-1}\Bigl(1 - \frac{k}{365}\Bigr),
\]
إذ يجب على كل شخص جديد أن يتجنب الأيام $k$ المأخوذة أصلا. ومن
أجل $n = 23$: $\P(D_{23}) \approx 0.493$ --- فتقاسم عيد ميلاد صار أرجح من عدمه أصلا.
والحدس الذي يفسر صغر $23$: بأخذ اللوغاريتمات، $-\ln
\P(D_n) \approx \sum_{k<n}\frac k{365} =
\frac{\binom n2}{365}$، ويعطي
$\binom{23}2 = 253$ المقدار $253/365 \approx 0.693 \approx \ln 2$. والمهم هو عدد \emph{الأزواج}، وهو ينمو
تربيعيا: فمسائل التصادم تعيش على السلّم $n \sim \sqrt{365}$ لا $n \sim
365$ --- ومنه
فمفارقة أعياد الميلاد جذر تربيعي متنكر.
\end{example}

\begin{example}[مونتي هول، ببايز]\label{ex:b2:proba:montyhall}
تختبئ جائزة وراء أحد أبواب ثلاثة، بانتظام. تختار الباب $1$؛
فيفتح المضيف، وهو يعرف مكان الجائزة، أحد البابين الآخرين، وهو دائما
فارغ (مع اختيار منتظم عندما يكون له خيار)، وليكن الباب $3$.
ولنجعل $B_i = {}$``الجائزة وراء الباب $i$'' و$A = {}$``يفتح المضيف
الباب $3$''. عندئذ $\pcond{B_1}{A} = \frac12$، $\pcond{B_2}{A} = 1$، $\pcond{B_3}{A} = 0$، ومنه ببايز (\cref{thm:b2:proba:bayes})،
\[
\P(B_2 \mid A)
= \frac{1\cdot\frac13}
{\frac12\cdot\frac13 + 1\cdot\frac13 + 0\cdot\frac13}
= \frac23 :
\]
فتبديل الأبواب يربح مرتين من ثلاث. ويحدد الحساب موضع الالتباس الشائع
بالضبط: فحركة المضيف \emph{مُخبِرة} (إذ لم يكن ليفتح الباب $2$
لو كانت الجائزة هناك)، وصيغة بايز هي أداة مسك الحسابات التي تحول هذا
اللاتناظر إلى $\frac23$. فالشرط على ``ما شوهد'' لا على ``ما هو صحيح''
هو فن الصيغة كله.
\end{example}

\begin{example}[رهانا الفارس دو ميريه]\label{ex:b2:proba:demere}
رهانان من القرن السابع عشر، يحسمهما الاستقلال. الرهان الأول: ظهور
ستة واحدة على الأقل في $4$ رميات لنرد،
\[
\P = 1 - \Bigl(\frac56\Bigr)^{\!4} \approx 0.518 > \frac12 .
\]
والرهان الثاني: ظهور ستة مضاعفة واحدة على الأقل في $24$ رمية
لنردين،
\[
\P = 1 - \Bigl(\frac{35}{36}\Bigr)^{\!24} \approx 0.491 <
\frac12 .
\]
وقد استدل دو ميريه بأن $24$ رمية بحظ $\frac1{36}$ ينبغي أن توافق $4$
رمية بحظ $\frac16$ (فالنسبة نفسها $\frac{24}{36} = \frac46$)؛ ويُقال إن فشل هذا التناسب
--- فاحتمالات الاتحادات لا تتحاكى خطيا --- هو الذي دفعه إلى مراسلة
باسكال، ومن ثم إلى ميلاد نظرية الاحتمالات. والمقارنة الصحيحة تمر عبر
اللوغاريتمات: فإن $n$ تجربة بحظ $p$ تنجح مرة واحدة على الأقل
باحتمال $1 - (1-p)^n \approx
1 - \eu^{-np}$، ومنه فإن الثابت الصادق هو $np$: وهنا $4\cdot\frac16 = \frac23$
مقابل $24\cdot\frac1{36} =
\frac23$ --- متساويان! فالرهانان لا يختلفان إلا في الرتبة الثانية
بدلالة $p$، وبما يكفي بالكاد لنقل أحدهما عبر خط الخمسين في
المائة: فالاحتمالات الصغيرة ميدان يحتاج فيه الحدس إلى الأسّي، لا إلى
المسطرة.
\end{example}

\begin{remark}[مغالطات شائعة في الشرط]
ثلاثة التباسات متكررة، وكلها ظاهرة في الأمثلة أعلاه. (1)
\emph{القلب}: يختلف $\pcond BA$ عن $\pcond AB$ بالعامل $\P(A)/\P(B)$ --- فاختبار
دقته $99\%$ على المرضى قد يترك مريضا نتيجته موجبة سليما شبه أكيد
عندما يكون المرض نادرا (\cref{exo:b2:proba:3})؛ وذكر $\pcond{\text{مريض}}{
\text{موجب}}$
في موضع يُقصد فيه $\pcond{\text{موجب}}{
\text{مريض}}$ هو مغالطة المعدل
القاعدي. (2) \emph{الشرط على \omterm{def:b2:proba:space}{الحدث} الخاطئ}: ففي مونتي هول، \omterm{def:b2:proba:space}{حدث} الشرط
الصحيح هو ``فتح المضيف الباب $3$''، لا ``الجائزة ليست وراء الباب
$3$''؛ فالحدثان يحملان معلومتين مختلفتين، وعلى الفرق يتوقف $\frac23$
كله. (3) \emph{المنفصل مقابل المستقل}: فالحدثان المنفصلان ذوا
الاحتمال الموجب ليسا مستقلين أبدا ($\P(A\cap B) = 0 \neq \P(A)\P(B)$) --- إذ إن الاستقلال توافق
في المعلومات، لا غياب للتراكب.
\end{remark}

\begin{definition}[الاستقلال]\label{def:b2:proba:independence}
يكون الحدثان $A$ و$B$ \emph{مستقلين}\index{أحداث مستقلة} إذا
كان $\P(A \cap B) =
\P(A)\P(B)$. وتكون عائلة $(A_i)_{i \in I}$ من \omterm{def:b2:proba:space}{الأحداث} \emph{مستقلة (فيما
بينها)} إذا كان من أجل كل جزء منته $J
\subseteq I$،
\[
\P\Bigl(\bigcap_{i \in J} A_i\Bigr)
= \prod_{i \in J} \P(A_i) .
\]
\end{definition}

\begin{remark}
الاستقلال المتبادل أقوى تماما من الاستقلال مثنى مثنى: فمع رميتين
لقطعة نقود متزنة، تكون \omterm{def:b2:proba:space}{الأحداث} ``الأولى صورة'' و``الثانية صورة''
و``الرميتان متفقتان'' مستقلة مثنى مثنى (فاحتمال تقاطع كل زوج $\frac14 =
\frac12\cdot\frac12$)،
ومع ذلك فإن احتمال التقاطع الثلاثي $\frac14 \neq \frac18$. ولاحظ أيضا أنه إذا كان
$A, B$ مستقلين، فكذلك $A, B^c$ (بالحساب: $\P(A \cap B^c) = \P(A) - \P(A\cap B) =
\P(A)(1 - \P(B))$)، ومنه أيضا $A^c, B^c$.
\end{remark}

\begin{example}[الاستقلال يُقرأ من بنية جدائية]\label{ex:b2:proba:productindep}
نرمي نردين متزنين: $\Omega = \intint16^2$ بأوزان منتظمة. ولنجعل $A = {}$``النرد
الأول زوجي'' و$B = {}$``النرد الثاني لا يقل عن $5$''. وبالعد:
$\abs A = 3\cdot6 = 18$، $\abs B = 6\cdot2 = 12$، $\abs{A\cap B} = 3\cdot2 = 6$، ومنه
\[
\P(A\cap B) = \frac6{36} = \frac{18}{36}\cdot\frac{12}{36} =
\P(A)\,\P(B) :
\]
فهما مستقلان، والآلية ظاهرة --- \omterm{def:b2:proba:space}{فالحدث} $A$ يقيد الإحداثية الأولى
وحدها، \omterm{def:b2:proba:space}{والحدث} $B$ الثانية وحدها، ويجعل القياس المنتظم على مجموعة
جدائية العدود الإحداثية تتضارب. وكل ادعاء من نوع ``\omterm{def:b2:proba:space}{الأحداث} المتعلقة
بمجموعات منفصلة من الرميات مستقلة'' (المستعمل بكثافة في مسألة نهاية
الأسبوع) هو هذا الحساب، مرتديا مزيدا من الأدلة.
\end{example}

\begin{example}[تحليل الخطوة الأولى]\label{ex:b2:proba:firststep}
من أجل النموذج الهندسي في \cref{ex:b2:proba:geometric}، ما احتمال $u$ أن تقع أول صورة
في رتبة \emph{زوجية}؟ نشترط على الرمية الأولى: فباحتمال $p$ تكون
الرتبة $1$ (فردية)؛ وباحتمال $q = 1 - p$ تبدأ اللعبة من جديد مع انقلاب
جميع الزوجيات، ومنه
\[
u = p\cdot0 + q\,(1 - u)
\qquad\Longrightarrow\qquad
u = \frac{q}{1 + q} .
\]
سطر واحد، دون أي متسلسلة --- وهو يوافق الجمع المباشر للمتسلسلة
\cref{exo:b2:proba:9}، الذي يعطي $1 - u =
\frac1{1+q}$. وهذه التقنية ``ذات الخطوة الأولى'' (اشترط
على التجربة الأولى، واعرف نسخة مزاحة من المسألة) هي الصورة الاحتمالية
للتراجع، وهي المحرك وراء معادلات مدة اللعبة في \cref{exo:b2:proba:6} وحسابات أول
عبور في مسألة نهاية الأسبوع.
\end{example}

\section{مبرهنة بوريل--كانتيلي}

\begin{definition}[النهاية العليا للأحداث]\label{def:b2:proba:limsup}
من أجل متتالية $(A_n)$ من \omterm{def:b2:proba:space}{الأحداث}، يكون \omterm{def:b2:proba:space}{الحدث} \[
\limsup_n A_n
= \bigcap_{N=0}^{\infty}\ \bigcup_{n \geq N} A_n
= \{\omega \in \Omega : \omega \in A_n
\text{ من أجل عدد
لانهائي من } n\}
\]
هو \omterm{def:b2:proba:space}{الحدث} ``يقع $A_n$ عددا لانهائيا من المرات''.
\end{definition}

\begin{example}[ترجمة ``عددا لانهائيا من المرات'' و``في النهاية'']\label{ex:b2:proba:limsuplang}
متمم $\limsup_nA_n$ هو، بحسب دي مورغان، \[
\Bigl(\bigcap_N\bigcup_{n\geq N}A_n\Bigr)^{\!c}
= \bigcup_N\bigcap_{n\geq N}A_n^c
= \{\omega : \omega \notin A_n \text{ من أجل كل }n\},
\] كبيرة
أي \omterm{def:b2:proba:space}{الحدث} ``\emph{في النهاية}، يفشل $A_n$'' (ويُكتب $\liminf_nA_n^c$). ومنه
فإن ``$A_n$ عددا لانهائيا من المرات'' و``$A_n^c$ في النهاية''
متتامان --- والحفاظ على وضوح هذا المعجم يمنع معظم حوادث المكمّمات.
وترجمات نموذجية من أجل رمي قطعة نقود: ``عدد لانهائي من الصور'' هو
$\limsup\{X_n = H\}$؛ و``عدد منته فقط من السلاسل ذات $100$ صورة'' هو متمم نهاية
عليا؛ و``التواتر الجاري يتقارب نحو $\frac12$'' هو $\bigcap_j\bigcup_N\bigcap_{n\geq N}\{\abs{\widehat p_n -
\tfrac12} < \tfrac1j\}$ --- وكلها
عمليات \omterm{def:b2:structures:countable}{قابلة للعد}، ومنه فكل هذه أحداث صادقة.
\end{example}

\begin{theorem}[بوريل--كانتيلي]\label{thm:b2:proba:borelcantelli}
\leavevmode
\begin{enumerate}
\item إذا كان $\sum_{n} \P(A_n) < \infty$، فإن $\P\bigl(\limsup_n A_n\bigr) = 0$.
\item إذا كانت \omterm{def:b2:proba:space}{الأحداث} $A_n$ مستقلة وكان $\sum_n \P(A_n) =
\infty$، فإن $\P\bigl(\limsup_n A_n\bigr) = 1$.
\end{enumerate}
\end{theorem}

\begin{proof}
\emph{1.} نضع $C_N = \bigcup_{n \geq N}A_n$؛ فالمتتالية $(C_N)$ متناقصة وتقاطعها $\limsup A_n$،
وبالجمعية التحتية \omterm{def:b2:structures:countable}{القابلة للعد} (\cref{cor:b2:proba:subadd})
\[
\P(C_N) \leq \sum_{n \geq N}\P(A_n)
\xrightarrow[N\to\infty]{} 0
\]
(وهو ذيل متسلسلة متقاربة). ويختم الاتصال الرتيب (\cref{thm:b2:proba:continuity}): $\P(\limsup A_n) =
\lim_N \P(C_N) = 0$.

\emph{2.} يكفي أن نبين $\P\bigl(\bigcup_{n\geq N}A_n\bigr)
= 1$ من أجل كل $N$: فإذا كانت أحداث
$B_N$ كلها ذات احتمال $1$، فإن
\[
\P\Bigl(\Bigl(\bigcap_NB_N\Bigr)^{\!c}\Bigr)
= \P\Bigl(\bigcup_NB_N^c\Bigr)
\leq \sum_N\P(B_N^c) = 0
\]
بالجمعية التحتية \omterm{def:b2:structures:countable}{القابلة للعد} (\cref{cor:b2:proba:subadd})، ومنه يبقى احتمال التقاطع
القابل للعد $\limsup A_n = \bigcap_N\bigcup_{n\geq N}A_n$ مساويا $1$. نثبّت $N$، وننظر من أجل $M > N$
في المتمم:
\[
\P\Bigl(\bigcap_{n=N}^{M} A_n^c\Bigr)
= \prod_{n=N}^{M}\bigl(1 - \P(A_n)\bigr)
\leq \prod_{n=N}^{M} e^{-\P(A_n)}
= \exp\Bigl(-\sum_{n=N}^M \P(A_n)\Bigr) ,
\]
باستعمال استقلال المتممات وحصر التحدب $1 -
x \leq e^{-x}$. وعندما $M \to \infty$ يؤول
الأسّ إلى $-\infty$ بحكم تباعد المتسلسلة، ومنه بالاتصال الرتيب
(للمتتالية المتناقصة) $\P\bigl(\bigcap_{n \geq N}A_n^c\bigr) = 0$، أي $\P\bigl(\bigcup_{n \geq N}A_n\bigr) = 1$.
\end{proof}

\begin{example}[سلاسل لانهائية من الصور]\label{ex:b2:proba:runs}
نرمي قطعة نقود متزنة إلى الأبد، ونجعل $A_n$ \omterm{def:b2:proba:space}{الحدث} ``الرميات $n,
n+1, \dots, n + k - 1$
كلها صور'' (أي سلسلة من $k$ صورة تبدأ عند الزمن $n$)، من أجل
$k$ ثابتة. \omterm{def:b2:proba:space}{والأحداث} $A_{jk}$ ($j =
1, 2, \dots$)، المتعلقة بكتل منفصلة من
الرميات، مستقلة، واحتمال كل منها $2^{-k}$، و$\sum_j 2^{-k} =
\infty$: ومنه، بحسب
بوريل--كانتيلي~2، يكون عدد لانهائي من الكتل كله صورا باحتمال $1$
--- أي إن \emph{كل} نمط ثابت يتكرر عددا لانهائيا من المرات، بشكل شبه
أكيد. وبالعكس، إذا تركنا طول السلسلة ينمو، فإن $B_n = {}$``تبدأ سلسلة
من $2\log_2 n$ صورة عند $n$'' له $\P(B_n) = n^{-2}$ قابلا للجمع، ومنه فإن عددا
منتهيا فقط من هذه السلاسل الطويلة يبدأ بشكل شبه أكيد: فبوريل--كانتيلي
تعاير بدقة \emph{كم يبلغ طول} أطول السلاسل.
\end{example}

\begin{example}[القرد اللانهائي، مقيسا]
يطبع قرد حروفا مستقلة منتظمة من أبجدية ذات $26$ حرفا. ونقطّع
المطبوع إلى كتل منفصلة من أربعة حروف؛ فتكون \omterm{def:b2:proba:space}{الأحداث} $A_j = {}$``تهجئ
الكتلة $j$ كلمة رياض'' مستقلة مع $\P(A_j) = 26^{-4}$، و$\sum_j\P(A_j) = \infty$: ومنه، بحسب
بوريل--كانتيلي~2، يكتب القرد ``رياض'' عددا لانهائيا من المرات بشكل
شبه أكيد --- والأمر نفسه يصح من أجل أي نص ثابت مهما يكن طوله، مع
تعديل الكتل. والحاشية الكمية تفرّغ المعجزة من هوائها: فالمقدار
$26^4 =
456\,976$، ومنه فإن أول ``رياض'' يستغرق نحو نصف مليون ضربة مفتاح وسطيا،
ومسرحية لشكسبير من $10^5$ حرفا تنتظر رتبة $26^{10^5}$ كتلة --- فالتأكد
شبه الأكيد نص عن الأفق $\infty$، لا عن أي أفق قد يبلغه قرد. فبوريل--كانتيلي
تصادق على النهاية؛ أما حجم الحدود فيروي القصة على السلالم البشرية.
\end{example}

\begin{remark}
في \cref{ex:b2:proba:runs} يكون \omterm{def:b2:proba:space}{الفضاء العيني} الأساسي (متتاليات الرميات اللانهائية)
غير قابل للعد، ومنه فإن المثال يعيش، بدقة الكلام، في إطار نظرية
القياس في السنة الثالثة؛ غير أن \emph{الحسابات} لا تستعمل إلا القواعد
المبرهنة في هذا الفصل، مطبقة على أحداث يحددها عدد منته من الرميات
وعلى توليفاتها \omterm{def:b2:structures:countable}{القابلة للعد}. وهذا هو الاصطلاح المعتاد في هذا المستوى:
تُذكر النظرية على الفضاءات \omterm{def:b2:structures:countable}{القابلة للعد}، وتُعالج أمثلة الألعاب
اللانهائية بالعدة نفسها.
\end{remark}

\begin{remark}[آفاق داخل هذا المجلد]
تستهلك آلة هذا الفصل بالجملة في الفصلين التاليين. فالدوال المميزة
تحول \omterm{def:b2:proba:space}{الأحداث} إلى متغيرات عشوائية، وتصير الجمعية $\sigma$ هي القابلية
للجمع التي تعرّف الأمل الرياضي (\cref{ch:b2:randomvar})؛ ومبرهنة بوريل--كانتيلي مع
حصر ذيلي قابل للجمع هي بالضبط كيف يُبرهن هناك على القانون القوي
للأعداد الكبيرة من أجل قطع النقود. وفي \cref{ch:b2:genfun}، يعود الاتصال الرتيب
في اللحظة الحاسمة: إذ \emph{يُعرَّف} احتمال انقراض مسار تفرع بأنه
النهاية الرتيبة $\lim\P(Z_n = 0)$، وتُحصَّل معادلة النقطة الثابتة التي يحققها
بالمرور إلى النهاية في تلك المتتالية المتزايدة --- فمبرهنة الكتاب
الأخيرة تقوم على مبرهنة هذا الفصل الأولى.
\end{remark}

\begin{remark}[طريقة: ثلاث طرائق إلى الاحتمال واحد]
تُبرهن النصوص شبه الأكيدة بثلاث روافع، بترتيب متزايد في القوة.
\emph{الاتصال الرتيب}: أظهر \omterm{def:b2:proba:space}{الحدث} اتحادا متزايدا (أو تقاطعا
متناقصا) لأحداث ذات أفق منته واحتمالات قابلة للحساب (\cref{ex:b2:proba:sixeventually}).
\emph{الاتحادات المعدومة}: اتحاد قابل للعد من أحداث معدومة الاحتمال
معدوم (بالجمعية التحتية \omterm{def:b2:structures:countable}{القابلة للعد})، ومنه يكفي قتل كل \omterm{def:b2:proba:space}{حدث} سيئ على
حدة --- وهكذا تتجمع عبارة ``من أجل كل $j$، وفي النهاية $\abs{\widehat p_n - p} < 1/j$''
في تقارب. \emph{بوريل--كانتيلي}: عندما يكون \omterm{def:b2:proba:space}{الحدث} نهاية عليا، اجمع
الاحتمالات؛ فالتقارب يقتله (دون حاجة إلى استقلال)، والتباعد مع
الاستقلال يصادق عليه. واختيار الرافعة الصحيحة هو عادة البرهان كله؛
وتشغّل مسألة نهاية الأسبوع الروافع الثلاث في حجة واحدة.
\end{remark}

\begin{remark}[أين يُستعمل هذا]
الاتصال الرتيب وبوريل--كانتيلي هما رافعتا كل نص ``شبه أكيد'': فهما
يقودان عوْد \omterm{pb:b2:proba:1}{السير العشوائي} في مسألة نهاية الأسبوع لهذا الفصل، والشق
شبه الأكيد من قانون الأعداد الكبيرة (\cref{ch:b2:randomvar})، وتحليل انقراض مسارات
التفرع (\cref{ch:b2:genfun}). ويعيد مجلد السنة الثالثة بناء النظرية على جبور
$\sigma$ وتكامل لوبيغ، حيث تصير الفضاءات العينية غير \omterm{def:b2:structures:countable}{القابلة للعد}
المستعملة هنا بصورة غير رسمية دقيقةً تماما.
\end{remark}

\section{تمارين}

\begin{exercise}[$\star$]\label{exo:b2:proba:1}
تحوي جرة $n$ كرة مرقمة. وتُسحب الكرات واحدة واحدة دون إرجاع.
احسب احتمال أن تُسحب الكرة رقم $1$ قبل الكرة رقم $2$. وعمّم:
احتمال أن تُسحب الكرة $1$ أولا من بين الكرات $1, \dots, k$.
\end{exercise}

\begin{exercise}[$\star$]\label{exo:b2:proba:2}
بيّن أن الأوزان $p_k = \frac{1}{k(k+1)}$ تعرّف على $\Omega = \N^*$ قياسا احتماليا، واحسب $\P(2\N^*)$
(النتائج الزوجية) على صورة متسلسلة؛ وبيّن أنه يساوي $1 - \ln 2$.
\emph{(تلسكب $\frac{1}{2j(2j+1)} = \frac{1}{2j} -
\frac{1}{2j+1}$ واستعمل المتسلسلة التوافقية المتناوبة، \cref{ch:b2:series}.)}
\end{exercise}

\begin{exercise}[$\star$]\label{exo:b2:proba:3}
(الإيجابيات الكاذبة) يصيب مرض شخصا من كل $10\,000$. ويكشفه اختبار
باحتمال $0.99$ عند المرضى، ويعطي إيجابية كاذبة باحتمال $0.01$ عند
الأصحاء. احسب احتمال أن يكون المرء مريضا علما بأن الاختبار موجب،
وعلّق.
\end{exercise}

\begin{exercise}[$\star\star$]\label{exo:b2:proba:4}
لتكن $A_1, \dots, A_n$ أحداثا. برهن على صيغة الاحتواء والاستبعاد
\[
\P\Bigl(\bigcup_{i=1}^n A_i\Bigr)
= \sum_{\emptyset \neq J \subseteq \{1,\dots,n\}}
(-1)^{\abs J + 1}\,\P\Bigl(\bigcap_{i \in J}A_i\Bigr)
\]
بمكاملة المتطابقة $1 - \prod_{i=1}^n(1 -
\mathbf{1}_{A_i}) = \mathbf{1}_{\bigcup A_i}$ على $\Omega$ (أي بالجمع موزونا بالمقدار
$\P(\{\omega\})$).
\end{exercise}

\begin{exercise}[$\star\star$]\label{exo:b2:proba:5}
(مسألة التطابقات، بالاحتواء والاستبعاد) تُوضع $n$ رسالة بانتظام
وعشوائيا في $n$ ظرفا، رسالة في كل ظرف. باستعمال \cref{exo:b2:proba:4}، بيّن أن
احتمال ألا يقع \emph{أي} تطابق صحيح هو $\sum_{k=0}^n \frac{(-1)^k}{k!} \to e^{-1}$، واستنتج احتمال وقوع
تطابق واحد بالضبط.
\end{exercise}

\begin{exercise}[$\star\star$]\label{exo:b2:proba:6}
تُرمى قطعة نقود منحازة (احتمال الصورة فيها $p \in \intoo{0}{1}$) حتى تظهر صورتان
متتاليتان. وليكن $q_n$ احتمال أن تدوم اللعبة أكثر من $n$ رمية.
بيّن، بالشرط على الرمية (أو الرميات) الأولى، أن $q_n = (1-p)\,q_{n-1} + p(1-p)\,q_{n-2}$ من أجل $n \geq 2$،
واستنتج أن اللعبة تنتهي باحتمال $1$. \emph{(بيّن $q_n \to 0$ بالمقارنة
مع متتالية هندسية: فجذرا المعادلة المميزة كلاهما في $\intoo{0}{1}$ بالقيمة
المطلقة.)}
\end{exercise}

\begin{exercise}[$\star\star\star$]\label{exo:b2:proba:7}
(الأرقام القياسية) نسحب متتالية لانهائية من الترتيبات المنتظمة
المستقلة، بالمعنى التوفيقي التالي: من أجل كل $n$، يكون الترتيب
النسبي لأول $n$ سحبة منتظما بين الإمكانات $n!$، ويكون
$R_n = {}$``السحبة رقم $n$ رقم قياسي (أكبر من جميع سابقاتها)''.
وبالتسليم بأن \omterm{def:b2:proba:space}{الأحداث} $R_n$ مستقلة مع $\P(R_n) = 1/n$ (وبرهن على الأقل على
هذه المساواة الأخيرة بالتناظر)، بيّن باستعمال بوريل--كانتيلي أن عددا
لانهائيا من الأرقام القياسية يقع بشكل شبه أكيد، لكن أن الأرقام
القياسية في لحظتين \emph{متتاليتين} $n, n+1$ تقع عددا لانهائيا من
المرات باحتمال --- احسب $\sum_n \P(R_n \cap R_{n+1})$ واستنتج ماذا تعطي بوريل--كانتيلي~1.
\end{exercise}

\begin{exercise}[$\star\star\star$]\label{exo:b2:proba:8}
(بنكهة كوخن--ستون، بصيغة أسهل) لتكن $(A_n)$ أحداثا مستقلة مع $\P(A_n) = \frac{1}{n+1}$.
بيّن أن $\P(\limsup A_n) = 1$، رغم أن $\P(A_n) \to 0$: أي ``نادرة فرادى، مؤكدة جماعة''.
وبالعكس، أظهر متتالية من \omterm{def:b2:proba:space}{الأحداث} (المتعلقة) مع $\sum\P(A_n) = \infty$ و$\P(\limsup A_n) = 0$، مما
يبين أنه لا يمكن إسقاط الاستقلال في بوريل--كانتيلي~2.
\end{exercise}

\begin{exercise}[$\star$]\label{exo:b2:proba:9}
تُرمى قطعة نقود احتمال الصورة فيها $p \in \intoo01$ حتى أول صورة. احسب احتمال
أن يقع ذلك في رتبة فردية، وقوّمه من أجل قطعة متزنة.
\end{exercise}

\begin{exercise}[$\star\star$]\label{exo:b2:proba:10}
لتكن $(A_n)_{n\geq1}$ أحداثا مستقلة مع $\P(A_n) =
p_n < 1$. بيّن أن
\[
\P\Bigl(\bigcap_{n\geq1}A_n^c\Bigr)
= \prod_{n\geq1}(1 - p_n)
:= \lim_{N\to\infty}\prod_{n=1}^N(1 - p_n),
\]
وأن هذه النهاية تساوي $> 0$ إذا وفقط إذا كان $\sum p_n <
\infty$. ووفّق ذلك مع
بوريل--كانتيلي: فعندما $\sum p_n =
\infty$، لا يقع بشكل شبه أكيد \omterm{def:b2:proba:space}{حدث} $A_n$ واحد
فحسب --- بل يقع عدد لانهائي منها.
\end{exercise}

\begin{exercise}[$\star\star$]\label{exo:b2:proba:11}
(علبة كبريت بناخ) يحتفظ مدخن بعلبة فيها $n$ عود كبريت في كل جيب،
ويمد يده في كل مرة إلى جيب مختار بانتظام عشوائيا. وعندما يجد علبة
فارغة أول مرة، ما احتمال أن تحوي العلبة الأخرى $k$ عودا بالضبط؟
بيّن أن الجواب هو $\binom{2n-k}{n}2^{-(2n-k)}$ وتحقق من أن مجموع هذه الاحتمالات $1$
من أجل $n = 1$.
\end{exercise}

\begin{exercise}[$\star\star\star$]\label{exo:b2:proba:12}
(الجمعية $\sigma$ بديهية حقيقية)
(1) بيّن أنه لا يوجد \omterm{def:b2:proba:space}{قياس احتمالي} على $(\N,
\mathcal P(\N))$ يعطي جميع المجموعات
الأحادية الوزن نفسه.
(2) من أجل $A \subseteq \N^*$، نضع $d(A) =
\lim_n\frac{\abs{A\cap\intint1n}}{n}$ عندما توجد النهاية (وهي \emph{الكثافة
الطبيعية}). بيّن أن $d$ جمعي على الأزواج التي توجد فيها الكثافات
الثلاث، وأنه يعطي كل مجموعة أحادية الكثافة $0$ ويعطي $\N^*$
الكثافة $1$ --- واستنتج أن $d$ ليس $\sigma$-جمعيا.
(3) أظهر مجموعة لا كثافة لها. \emph{(بمناوبة الكتل $\intint{2^{2k}}{2^{2k+1}-1}$ داخلا
وخارجا.)}
\end{exercise}

\section{مسألة: السير العشوائي البسيط على $\Z$ عوّاد}

\begin{omfigure}
\begin{tikzpicture}[scale=0.42, y=0.9cm]
  \draw[->] (0,0) -- (25,0) node[below] {$n$};
  \draw[->] (0,-2.6) -- (0,2.9) node[left] {$S_n$};
  \draw[omDef, thick] plot coordinates {(0,0) (1,1) (2,0)
    (3,-1) (4,0) (5,1) (6,2) (7,1) (8,2) (9,1) (10,0) (11,-1)
    (12,-2) (13,-1) (14,0) (15,-1) (16,0) (17,1) (18,2) (19,1)
    (20,0) (21,1) (22,0) (23,1) (24,2)};
  \foreach \t in {2,4,10,14,16,20,22}
    \fill[omThm] (\t,0) circle (5pt);
  \fill (0,0) circle (5pt);
\end{tikzpicture}

{\small أربع وعشرون خطوة من \omterm{pb:b2:proba:1}{سير عشوائي} بسيط؛ والنقط الحمراء تعلّم
العودات إلى المبدأ. وتبين المسألة أن هذه النقط، باحتمال $1$، لا
تتوقف عن الظهور أبدا --- ومع ذلك فإن زمن الانتظار بينها متوسطه
متباعد.}
\end{omfigure}

\begin{problem}[{مسألة نهاية الأسبوع --- مبرهنة بوليا في العوْد على
$\Z$، مع مسألة الاقتراع ونكهة قوس الجيب على الطريق}]\label{pb:b2:proba:1}
نرمي قطعة نقود متزنة إلى الأبد؛ وليكن $X_i = \pm1$ الخطوة رقم $i$
وليكن $S_n = X_1 + \dots + X_n$ \emph{السير العشوائي البسيط}\index{سير عشوائي} على
$\Z$، $S_0 = 0$. وكما في \cref{ex:b2:proba:runs}، تُحدَّد جميع \omterm{def:b2:proba:space}{الأحداث} أدناه بعدد
منته من الرميات أو تكون توليفات \omterm{def:b2:structures:countable}{قابلة للعد} من هذه \omterm{def:b2:proba:space}{الأحداث}، ويكون
استقلال \omterm{def:b2:proba:space}{الأحداث} المتعلقة بكتل منفصلة من الرميات جزءا من النموذج.
ونكتب $u_n =
\P(S_{2n} = 0)$ و$N_n(k)$ لعدد المسارات $\pm1$ ذات الطول $n$ من
$0$ إلى $k$.

\medskip
\noindent\textbf{الجزء الأول --- عدّ المسارات.}
\begin{enumerate}
  \item بيّن أن $N_n(k) = \binom{n}{(n+k)/2}$ عندما يكون $n + k$ زوجيا و$\abs k \leq n$، وأن $0$
        فيما عدا ذلك؛ واستنتج $\P(S_n = k) = N_n(k)\,2^{-n}$. ولماذا يكون لكل مسار مفرد ذي
        الطول $n$ الاحتمال نفسه؟
  \item بيّن $S_{2n+1} \neq 0$ و$u_n =
        \binom{2n}{n}4^{-n}$، واحسب $u_1, u_2, u_3$.
  \item برهن على $u_n = \frac{2n-1}{2n}\,u_{n-1}$؛ واستنتج أن $(u_n)$ يتناقص نحو $0$،
        واستنتج من \cref{ex:b2:comparison:centralbinomial} أن
        \[
        u_n \sim \frac{1}{\sqrt{\pi n}},
        \qquad\text{ومنه}\qquad
        \sum_n u_n = \infty .
        \]
  \item (مبدأ الانعكاس) من أجل $k \geq 1$، بيّن أن مسارات الطول $n$
        من $1$ إلى $k$ التي تلامس $0$ في تقابل مع
        المسارات من $-1$ إلى $k$؛ واستنتج أن عدد المسارات من
        $0$ إلى $k$ التي تبقى $> 0$ بعد الزمن $0$ هو
        $N_{n-1}(k-1) -
        N_{n-1}(k+1)$.
  \item (مبرهنة الاقتراع) استنتج
        \[
        \P\bigl(S_1 > 0, \dots, S_{n-1} > 0 \bigm| S_n =
        k\bigr) = \frac kn \qquad (k \geq 1) :
        \]
        أي إنه في فرز يتقدم فيه الفائز بفارق $k$ من أصل $n$
        ورقة، يكون احتمال أن يتقدم الفائز طوال الفرز هو $k/n$.
        وتحقق يدويا من أجل $n = 3$، $k =
        1$.
\end{enumerate}

\medskip
\noindent\textbf{الجزء الثاني --- العودة إلى المبدأ.}
\begin{enumerate}[resume]
  \item برهن على المتطابقة المفتاحية
        \[
        \P(S_1 \neq 0,\ S_2 \neq 0,\ \dots,\ S_{2n} \neq 0) =
        u_n
        \]
        \emph{(اشترط على الخطوة الأولى، واجمع عدود السؤال 4 على
        نقطة الوصول، وتلسكب؛ واختم بالمقدار $2\binom{2n-1}{n} = \binom{2n}{n}$)}.
  \item استنتج من الاتصال الرتيب (\cref{thm:b2:proba:continuity}) أن السير يعود إلى
        $0$ مرة واحدة على الأقل باحتمال $1$، وأن
        $f_n := \P(\text{أول عودة عند الزمن }2n)$
        يحقق
        \[
        f_n = u_{n-1} - u_n = \frac{u_n}{2n-1},
        \qquad \sum_{n\geq1}f_n = 1 .
        \]
  \item بيّن أن $\sum_n 2n\,f_n = \infty$: فالعودة مؤكدة، لكن المتسلسلة التي كانت
        ستحسب متوسط زمن الانتظار متباعدة (وبمفردات \cref{ch:b2:randomvar}، يكون
        لزمن العودة أمل رياضي لانهائي).
  \item برهن على أنه من أجل كل $k \geq 1$، $\P(\text{ما لا يقل عن }
        k\text{
        عودة إلى }0) = 1$ \emph{(بالتفكيك على أزمنة العودات $k$
        الأولى: فكتل الرميات الموافقة منفصلة، ومنه تتضارب الاحتمالات
        ويكون مجموعها $(\sum_nf_n)^k$)}؛ واختم بالاتصال الرتيب:
        \[
        \P(S_n = 0 \text{ من أجل عدد لانهائي من } n) = 1 :
        \]
        أي إن \omterm{pb:b2:proba:1}{السير العشوائي البسيط} على $\Z$ \emph{عوّاد}.
  \item بيّن أن السير يزور كل موقع $k \in \Z$ بشكل شبه أكيد، ومنه (بحكم
        العوْد، بعد إعادة الانطلاق عند أول زيارة) عددا لانهائيا من
        المرات. \emph{(إشارات الجولات المتتالية انطلاقا من $0$
        قطع نقود متزنة مستقلة؛ والجولة الموجبة تزور $1$.)}
\end{enumerate}

\medskip
\noindent\textbf{الجزء الثالث --- بوريل--كانتيلي والسير المنحاز.}
\begin{enumerate}[resume]
  \item تحقق \omterm{def:b2:proba:space}{الأحداث} $A_n = \{S_{2n} = 0\}$ العلاقة $\sum\P(A_n)
        = \infty$؛ اشرح لماذا \emph{لا}
        تنطبق عليها بوريل--كانتيلي~2، وماذا كانت ستعطي
        بوريل--كانتيلي~1 لو تقاربت المتسلسلة. (وهذه هي استراتيجية
        الجزء كله.)
  \item ولنجعل الآن للقطعة انحيازا $p \neq \frac12$، $q = 1 -
        p$. بيّن $\P(S_{2n} = 0) = \binom{2n}n(pq)^n =
        u_n\,(4pq)^n$ مع
        $4pq < 1$، واستنتج $\sum_n\P(S_{2n} = 0) < \infty$، واختم ببوريل--كانتيلي~1 أن السير
        المنحاز لا يعود إلى $0$ إلا عددا منتهيا من المرات، بشكل
        شبه أكيد.
  \item ومن أجل $p \neq \frac12$ أيضا: بيّن $\P(S_n = k) \leq
        \binom{n}{\floor{n/2}}\,(pq)^{n/2}\,(p/q)^{k/2}$ من أجل كل $k$ ثابتة،
        واستنتج أن كل موقع يُزار عددا منتهيا من المرات بشكل شبه أكيد،
        واستنتج $\abs{S_n} \to \infty$ بشكل شبه أكيد: أي إن السير المنحاز
        \emph{عابر}.
  \item عودا إلى القطعة المتزنة: باستعمال السؤال 6، احسب احتمال ألا
        تنتج $200$ رمية \emph{أي} تعادل ($S_n \neq 0$ من أجل $1 \leq n \leq 200$)،
        وعدديا $u_{100} \approx 0.056$. وعلّق على التلاشي البطيء $1/\sqrt{\pi n}$: فالتعادلات
        مؤكدة على المدى الطويل لكنها أندر مما يوحي به الحدس.
  \item (أول عبور) ليكن $T_1$ أول زمن يبلغ فيه السير $1$.
        باستعمال مبدأ الانعكاس من أجل القيمة العظمى $M_n = \max_{i\leq n}S_i$ (المبرهن
        في السؤال 16، وهو لا يتعلق بهذا السؤال)، أو مباشرة انطلاقا
        من السؤال 7 بالشرط على الخطوة الأولى، بيّن $\P(T_1 = 2n - 1) = f_n$؛ واستنتج
        $\P(T_1 <
        \infty) = 1$ بينما تتباعد متسلسلة متوسط الزمن $\sum(2n-1)f_n$.
\end{enumerate}

\medskip
\noindent\textbf{الجزء الرابع --- القيم العظمى، وآخر صفر، والتقدمات
الطويلة.}
\begin{enumerate}[resume]
  \item (الانعكاس من أجل القيمة العظمى) من أجل $k \geq 1$، برهن على
        \[
        \P(M_n \geq k) = 2\,\P(S_n > k) + \P(S_n = k)
        \]
        بعكس المسار بعد أول زيارة له للمستوى $k$.
  \item استنتج $\P(M_{2n} \geq 1) = 1 - u_n$، أي $\P(S_i \leq 0 \text{ من أجل كل } i \leq 2n) = u_n$: فاحتمال ألا
        يتقدم السير أبدا يساوي احتمال ألا يكون عند الصفر أبدا
        (السؤال 6) --- حدثان مختلفان، واحتمال واحد.
  \item (آخر صفر) ليكن $L_{2n} = \max\{k \leq 2n : S_k =
        0\}$ (زوجيا). بجمع السؤال 6 مع استقلال
        كتل الرميات المنفصلة، بيّن
        \[
        \P(L_{2n} = 2k) = u_k\,u_{n-k}
        \qquad (0 \leq k \leq n),
        \]
        واستنتج، دون أي حساب إضافي، المتطابقة الثنائية $\sum_{k=0}^n u_ku_{n-k} = 1$.
  \item بيّن أن قانون $L_{2n}$ متناظر ($\P(L =
        2k) = \P(L = 2n - 2k)$)، وبيّن باستعمال $u_j \sim
        1/\sqrt{\pi j}$
        أن قيمتيه الحديتين هما أرجح قيمه. وجدول من أجل $n = 5$:
        $\P(L_{10} = 0)
        = u_5 \approx 0.246$ مقابل $\P(L_{10} = 4) = u_2u_3
        \approx 0.117$. وفسّر: ففي لعبة متزنة طويلة، يميل آخر
        تعادل إلى أن يكون مبكرا جدا أو متأخرا جدا --- فالتقدمات
        الطويلة هي القاعدة لا الاستثناء.
  \item اجمع الأسئلة من 16 إلى 19 في فقرة عن صورة التقلبات في السير
        المتزن: السلّم الانتشاري الذي يوحي به السؤال 3، وتأكد العودة
        مقابل تباعد متوسط زمن الانتظار، وثبات التقدمات بنكهة قوس
        الجيب.
\end{enumerate}

\medskip
\noindent\textbf{الجزء الخامس --- متطابقة التجدد ومبرهنة بوليا.}
\begin{enumerate}[resume]
  \item برهن، بتجزئة $\{S_{2n} = 0\}$ على زمن أول عودة، على \emph{متطابقة
        التجدد}
        \[
        u_n = \sum_{k=1}^n f_k\,u_{n-k} \quad (n \geq 1),
        \qquad\text{ومنه}\qquad
        U(x)\bigl(1 - F(x)\bigr) = 1 \quad (0 \leq x < 1),
        \]
        حيث $U(x) = \sum_{n\geq0}u_nx^n$ و$F(x) =
        \sum_{n\geq1}f_nx^n$ (وبرر أنصاف الأقطار وجداء المتسلسلتين
        مع \cref{ch:b2:powerseries}).
  \item استنتج \emph{ثنائية العوْد}: بجعل $x
        \to 1^-$ (بالنهايات الرتيبة
        للمتسلسلات ذات المعاملات الموجبة)،
        \[
        \sum_n u_n = \infty \iff \sum_n f_n = 1 ,
        \]
        وتحقق منها بمقابلتها بالأسئلة 3 و7 (السير المتزن) و12 (السير
        المنحاز).
  \item (البعد $2$) يأخذ السير البسيط على $\Z^2$ الخطوات $(\pm1, 0)$،
        $(0, \pm1)$ بانتظام. بيّن أن الإحداثيتين المدارتين $U_n = X_n + Y_n$ و$V_n = X_n
        - Y_n$
        تؤديان سيرين متزنين \emph{مستقلين} على $\Z$، واستنتج
        \[
        \P\bigl(S^{(2)}_{2n} = (0,0)\bigr) = u_n^2 \sim
        \frac1{\pi n},
        \qquad \sum_n u_n^2 = \infty ,
        \]
        واختم بالسؤالين 21 و22 (اللذين ينتقل برهاناهما حرفيا) أن
        السير على $\Z^2$ عوّاد.
  \item (البعد $3$) من أجل السير البسيط على $\Z^3$، سلّم بالتقدير
        الموضعي $\P(S^{(3)}_{2n} = 0) \leq
        C\,n^{-3/2}$ (المبرهن بمبرهنة النهاية الموضعية في مجلد
        السنة الثالثة). استنتج من بوريل--كانتيلي~1 أن السير على
        $\Z^3$ عابر، واذكر النتيجة الكاملة: \emph{مبرهنة بوليا} ---
        \omterm{pb:b2:proba:1}{فالسير العشوائي البسيط} عوّاد في البعدين $1$ و$2$،
        وعابر في البعد $3$ وما فوقه.
  \item تركيب. اسرد الدور المضبوط الذي أدّاه كل مما يلي: عدّ المسارات
        والانعكاس؛ والاتصال الرتيب؛ واستقلال كتل الرميات المنفصلة؛
        وبوريل--كانتيلي~1؛ ومتطابقة التجدد. وأي حقيقة \omterm{def:b2:powerseries:analytic}{تحليلية} وحيدة
        ($u_n \sim 1/\sqrt{\pi n}$، ومنه $\sum u_n = \infty$ لكن $\sum u_n^2 = \infty$ و$\sum n^{-3/2} < \infty$) تقرر بين العوْد
        والعبور في كل بعد؟
\end{enumerate}
\end{problem}
